\chapter{Detalle de la batería de pruebas}
\label{ap:capa1}

Este apéndice deja el detalle operativo de la evaluación: qué se comprobó, con
qué resultado y qué parte quedó como revisión manual.

\section{Núcleo determinista y persistencia}

Se comprueba que los modelos respetan el contrato mínimo
\texttt{decide}, \texttt{update} y \texttt{get\_state}; que dos ejecuciones con
la misma semilla generan la misma traza; y que cada observación queda registrada
con el mismo identificador en PostgreSQL y Qdrant. La comprobación de
persistencia cruzada aparece en la Tabla~\ref{tab:joinable}.

\begin{table}[H]
\begin{center}
\small
\renewcommand{\arraystretch}{1.16}
\begin{tabular}{|>{\raggedright\arraybackslash}p{6.2cm}|>{\centering\arraybackslash}p{2.5cm}|>{\centering\arraybackslash}p{2.5cm}|} \hline
\textbf{Métrica de persistencia} & \textbf{CASO1} & \textbf{CASO2} \\ \hline
Observaciones de simulación generadas & $13$ & $9$ \\ \hline
Filas nuevas en \texttt{simulation\_observations} (PostgreSQL) & $13$ & $9$ \\ \hline
Memorias densas escritas (Voyage, 1024 dim.) & $13$ & $9$ \\ \hline
Memorias \textit{sparse} escritas (BM25 nativo) & $13$ & $9$ \\ \hline
\end{tabular}
\caption[Fidelidad de la persistencia]{Fidelidad de la persistencia. La igualdad
entre las observaciones generadas y los tres almacenes valida que PostgreSQL,
Qdrant denso y Qdrant \textit{sparse} contienen el mismo número de observaciones
bajo los mismos identificadores.}
\label{tab:joinable}
\end{center}
\end{table}

\section{Regresión automática}

La suite combina pruebas unitarias, integración backend y recorridos de
interfaz. La sesión real con backbone se mantiene como comprobación manual
porque depende de servicios vivos. Los bloques de regresión usados se muestran
en la Tabla~\ref{tab:suite}.

\begin{table}[H]
\begin{center}
\small
\setlength{\tabcolsep}{5pt}
\renewcommand{\arraystretch}{1.16}
\begin{tabular}{|>{\raggedright\arraybackslash}p{4.6cm}|>{\raggedright\arraybackslash}p{1.6cm}|>{\raggedright\arraybackslash}p{1.4cm}|>{\raggedright\arraybackslash}p{4.0cm}|} \hline
\textbf{Bloque} & \textbf{Estado} & \textbf{Casos} & \textbf{Qué demuestra} \\ \hline
Backend (\textit{pytest}) & PASA & 387/387 & Dominio, agentes, Orchestrator, persistencia, recuperación, carga de modelos, anclaje determinista y Reporter. \\ \hline
Determinismo del motor (misma semilla) & PASA & 1/1 & Dos ejecuciones con la misma semilla y configuración producen eventos idénticos. \\ \hline
E2E web (modo \textit{mock}) & PASA & 4/4 & Pipeline completo en interfaz sin servicios externos. \\ \hline
E2E web (sesión real y backbone) & Manual & 13 comprobaciones & Sesión completa y vistas del Knowledge Backbone; requieren el \textit{backend} en ejecución. \\ \hline
\end{tabular}
\caption[Suite de regresión y comprobaciones]{Suite de regresión y
comprobaciones usadas como base de la evaluación. El estado \emph{Manual} indica
pruebas ejecutadas sobre un entorno vivo, no ausencia de comprobación.}
\label{tab:suite}
\end{center}
\end{table}

Las comprobaciones visuales de legibilidad, navegación y Knowledge Backbone se
revisan sobre la aplicación en ejecución.

\section{Fidelidad factual del Reporter}

El script \texttt{faithfulness.py} compara el PDF con los datos de origen:
comida consumida, pasos, recompensa total y energía final. En una prueba
controlada, el informe fiel obtiene fidelidad $1{,}00$; un informe adversario
con errores plantados baja a $0{,}33$. También se comprueba que el PDF evaluado
sea el compilado real con \texttt{tectonic}, no el documento de respaldo.

\section[Escenarios controlados]{Validación de comportamiento con escenarios controlados}
\label{ap:capa3}

El banco incluye además escenarios controlados sobre modelos de la
\emph{ejecución de muestra} de la primera fase. Son distintos de CASO1/CASO2 y
se usan como \textit{golden scenarios}: la predicción esperada se conoce de
antemano y se comprueba contra la simulación
~\cite{wilson2019tenrules,palminteri2017falsification}. Los modelos evaluados
se enumeran en la Tabla~\ref{tab:modelos-sujeto}.

\begin{table}[H]
\begin{center}
\small
\renewcommand{\arraystretch}{1.16}
\begin{tabular}{|>{\raggedright\arraybackslash}p{2.8cm}|>{\raggedright\arraybackslash}p{3.6cm}|>{\raggedright\arraybackslash}p{5.0cm}|} \hline
\textbf{Paradigma} & \textbf{Modelo} & \textbf{Predicción falsable} \\ \hline
Forrajeo óptimo & Teorema del valor marginal (MVT) & Abandona el parche cuando la tasa marginal cae bajo la media del entorno~\cite{charnov1976mvt} \\ \hline
Forrajeo óptimo & Amplitud de dieta & Regla cero-uno: ignora la presa pobre si abunda la buena~\cite{macarthur1966dietbreadth} \\ \hline
Aprendizaje por refuerzo & Q-learning tabular & Curva de aprendizaje: el rendimiento mejora con la experiencia \\ \hline
Aprendizaje por refuerzo & \textit{Actor-critic} & Aprende; el error TD actúa como señal de recompensa-predicción \\ \hline
\end{tabular}
\caption[Modelos de los escenarios controlados]{Modelos de la ejecución de
muestra de la primera fase y la predicción falsable que debe reproducir cada uno.}
\label{tab:modelos-sujeto}
\end{center}
\end{table}

Como referencia se añaden dos modelos manuales: \texttt{GreedyForagerOracle}
(techo de eficiencia) y \texttt{RandomModel} (suelo). Los resultados se
presentan en la Tabla~\ref{tab:gs-resultados}.

\begin{table}[H]
\begin{center}
\scriptsize
\renewcommand{\arraystretch}{1.12}
\resizebox{\textwidth}{!}{%
\begin{tabular}{|>{\raggedright\arraybackslash}p{1.75cm}|>{\raggedright\arraybackslash}p{2.25cm}|>{\raggedright\arraybackslash}p{4.55cm}|>{\raggedright\arraybackslash}p{1.1cm}|} \hline
\textbf{ID} & \textbf{Modelo} & \textbf{Métrica observada} & \textbf{Est.} \\ \hline
GS-OFT-1 & MVT & Residencia máx.\ 12 pasos, 8 abandonos con reinicio & PASA \\ \hline
GS-OFT-2 & MVT & Residencia media al partir $2{,}14 \rightarrow 2{,}73$ al subir el coste de viaje & PASA \\ \hline
GS-OFT-3 & Diet-breadth & Fracción de dieta $\{1\}$: $0{,}65$ (buena abundante) vs.\ $0{,}00$ (escasa) & PASA \\ \hline
GS-RL-1 & Q-learning & Tasa de recompensa $+0{,}46$; $\varepsilon$ decae $0{,}99\rightarrow 0{,}05$ & PASA \\ \hline
GS-RL-1 & Actor-critic & Tasa de recompensa $+0{,}01$; $\beta$ asciende $3\rightarrow 15$ & FALLA \\ \hline
GS-RL-1 & MVT (control) & Tasa de recompensa $+0{,}00$ (plano) & PASA \\ \hline
GS-RL-1 & Diet-breadth (control) & Tasa de recompensa $-0{,}01$ (plano) & PASA \\ \hline
\end{tabular}
}
\caption[Resultados de los escenarios controlados]{Resultados de los escenarios
controlados; seis de las siete predicciones se reproducen.}
\label{tab:gs-resultados}
\end{center}
\end{table}
